\documentclass[10pt,a4paper]{amsart}
\usepackage[margin=19mm]{geometry}
\usepackage{amsmath,amssymb,mathtools}
\usepackage[scr=rsfs,cal=euler]{mathalfa}
\usepackage{xfrac}
\usepackage[unicode,hidelinks,bookmarks=false]{hyperref}
\newcommand{\ie}{i.e.\ }
\newcommand{\cf}{cf.\ }
\newcommand{\resp}{respectively\ }
\newcommand{\Subsection}{\subsection}
\providecommand{\nobreakdash}{\nobreak\hbox{-}}
\setlength{\emergencystretch}{2em}
\multlinegap0pt
\usepackage{bm}
\usepackage{bbm}
\usepackage{dsfont}
\usepackage{xspace}
\usepackage{cancel}
\usepackage{mathtools}
\usepackage{amsthm}
\makeatletter
\@ifundefined{theorem}{\newtheorem{theorem}{Theorem}}{}
\@ifundefined{corollary}{\newtheorem{corollary}[theorem]{Corollary}}{}
\makeatother
\makeatletter
\newcommand{\Char}{\operatorname{char}}
\newcommand{\Jac}{\operatorname{Jac}}
\newcommand{\Z}{\mathbb{Z}}
\expandafter\ifx\csname introtheorem\endcsname\relax
\newenvironment{introtheorem}[1]{
	\def\intr@thmname{Theorem #1}
	\begin{intr@thm}}
	{\end{intr@thm}}
\fi
\expandafter\ifx\csname labelledconj\endcsname\relax
\newenvironment{labelledconj}[3][]
\fi
\makeatother
\title[Iwasawa Theory of graphs and their duals]{Iwasawa Theory of graphs and their duals}
\author{Debanjana Kundu, Katharina M\"uller}
\date{Source: arXiv:2410.11704v2 (CC BY 4.0)}
\begin{document}
\maketitle

\noindent\emph{Source and licence.} Iwasawa Theory of graphs and their duals. Version arXiv:2410.11704v2 (CC BY 4.0), distributed under the Creative Commons Attribution 4.0 International licence, as recorded in the arXiv licence metadata for this exact version and shown on the abstract page https://arxiv.org/abs/2410.11704v2. Full rights notice: licenses/arxiv-2410.11704-CC-BY-4.0.txt.\par\smallskip

\noindent\textit{Editorial scope.} Counted unit: the Corollary whose source label is \texttt{None} in the article (source0.tex lines 2504--2508, source label \texttt{None}), delivered together with the article's own complete original proof of it (source0.tex lines 2510--2550, one \texttt{proof} environment of 41 lines). No mathematics is written, repaired or re-derived: statement and proof are the article's own text line for line. Required context retained uncounted: thm:planar (lines 698--703); cor:app (lines 2464--2469). Every definition, display and label of those lines is retained unchanged. Omitted from this package and not redistributed: 1--697, 704--2463, 2470--2503, 2509--2509, 2551--2555. Nothing inside a retained excerpt is omitted or altered: every retained body is delivered exactly as the article's lines are written, so no omission receipt of any kind accompanies this package. Citations: the retained excerpts carry no citation command at all, so no bibliography entry of the article is transcribed and no reference is redistributed. Reference adaptation: none was needed and none was made; every reference inside the delivered unit resolves inside it, so no unresolved reference or citation is produced. Printed slips kept literal: no printed word, symbol, hypothesis or number of the counted statement or of its proof was altered. Layout and class: the article's own class is \texttt{amsart} with packages including amssymb, stmaryrd, amsmath, amscd, amsthm, amsbsy, geometry, commath, longtable, comment, enumerate, xy, hyperref, mathrsfs; the wrapper is the lane's pinned \texttt{amsart} editorial wrapper at 10pt on A4, which restates the theorem-like environments the retained text needs and adds the article's own macro definitions that the retained text needs (\texttt{\textbackslash Char}, \texttt{\textbackslash Jac}, \texttt{\textbackslash Z}), with extra wrapper packages (bm, bbm, dsfont, xspace, cancel, mathtools). The delivered numbering is the unit's own. Assets: no retained span contains a figure environment, a graphics-inclusion command, a PDF-page inclusion, a table or a TikZ picture; no image file is copied into this package, the package declares no asset dependency and no image file is redistributed with it. Licence: version arXiv:2410.11704v2 of this article is distributed under the Creative Commons Attribution 4.0 International licence, as recorded in the arXiv licence metadata for this exact version and shown on the abstract page https://arxiv.org/abs/2410.11704v2. A full rights notice is delivered at licenses/arxiv-2410.11704-CC-BY-4.0.txt. Characters: every retained excerpt is pure ASCII, so the delivered engine, XeLaTeX with its default Latin Modern text font, reports no missing character.
\begin{theorem}
\label{thm:planar}
Fix a finite group $G$.
Let $X$ be a finite undirected planar graph and $Y$ be the undirected derived graph of $X$ of a single\footnote{By this we mean that exactly one of the edges has a non-trivial voltage assignment.} \emph{unramified} voltage assignment.
Then $Y$ is planar.
\end{theorem}

\begin{corollary}
\label{cor:app}
Assume that $(X_n)_{n\in \mathbb{N}}$ is an \emph{unramified} $\Z_p$-tower of planar graphs.
Assume that for each $n$, $X_n^\vee/X^\vee$ is a branched covering. 
Then $T\mid \Char_\Lambda({\Jac}_\Lambda(X_\infty))$.
\end{corollary}

\begin{corollary}
{Let $X$ be a finite connected graph and let $\alpha$ be a single voltage assignment.
Let $(X_n)_{n\in \mathbb{N}}$ be the unramified $\Z_p$-tower obtained from $\alpha$.
Then $T\mid \Char_\Lambda(\Jac_\Lambda(X_\infty))$.}
\end{corollary}

\begin{proof}
{Let $e_0$ be the edge of $X$ having a non-trivial voltage assignment.
We fix an embedding of $X$ such that $e_0$ lies on the outer face.}

{By Theorem~\ref{thm:planar} all graphs $X_n$ are planar.
For each $X_n$ we fix an embedding in the plane.
Let $X_n^\vee $ be the corresponding dual graphs.
If we can show that $X_n^\vee/X^\vee $ is a branched covering, the claim will follow from Corollary~\ref{cor:app}. }

{Let $X'$ be the graph obtained from $X$ by deleting the edge $e_0$.
Then $X_n$ consists of $p^n$ copies of $X'$ (the different sheets) and the additional edges that restrict to $e_0$.
We draw $X_n$ as $p^n$-copies of $X'$ next to each other and then add the additional edges.
The graph $X_n$ has three different types of faces: 
\begin{enumerate}
    \item faces that lie completely in one sheet of $X_n$.
    \item one outer face
    \item one face that is neither of the first two.
\end{enumerate}
Let
\[
\pi_n\colon X_n \longrightarrow X
\]
be the covering of graphs.
This induces a surjective map of graphs
\[
\pi^\vee_n\colon X_n^\vee \longrightarrow X^\vee.
\]}
{Consider a face $F_{{n}}$ of $X_n$ of type (1), and fix a path $\mathsf{P}_{F_n}$ around its boundary counterclockwise.
Note that then $\pi_n(\mathsf{P}_{F_{{n}}})=\mathsf{P}_{F'}$ where $F'$ is a face of $X'$ (and can also be viewed as a face of $X$).
Every edge of $F_n$ has a corresponding edge in $X_n^\vee$ by construction of the dual graph.
Note that each edge starting or ending in $F_{{n}}$ corresponds to (projects to) exactly one edge in $X^\vee$.
In particular, the edges of a fixed face $F_{n}$ of type (1) and edges of $F'$ are in one to one correspondence.
We further note that the face $F'$ is a vertex in $X^\vee$.}

%\textcolor{blue}{Is this somehow saying that there is vertex in $X^\vee$ corresponding to the face $F'$?}\KM{exactly. Every face of $X'$ (except for the outer one is also a face of $X$}
{Now, consider a face $F_{n}$ of $X_n$ type (2).
Observe that $\pi_n(\mathsf{P}_{F_{n}})=p^n \mathsf{P}_O$, where $O$ is the outer face of $X$.
If $F_n$ is a face of type (3), then $\pi_n(\mathsf{P}_{F_n})=p^n \mathsf{P}_{F_0}$, where $F_0$ is the face of $X$ having $e_0$ on its boundary but is not the outer face. 
Thus there is a $p^n$-to-1 correspondence between edges starting/ending at $F_{n}$ and edges starting and ending at $O$.
Thus, $\pi^\vee_n$ is indeed a branched covering.}
\end{proof}

\end{document}
